\documentclass[11pt]{article}

\usepackage[margin=1in]{geometry}
\usepackage[T1]{fontenc}
\usepackage[utf8]{inputenc}
\usepackage{lmodern}
\usepackage{amsmath,amssymb,amsthm,mathtools,bm}
\usepackage{booktabs}
\usepackage{hyperref}
\usepackage{enumitem}
\usepackage{xcolor}
\usepackage{microtype}

\hypersetup{
  colorlinks=true, linkcolor=blue, urlcolor=blue, citecolor=blue,
  pdftitle={Diagnostic Metrics for the Grokking Transition},
}

\newcommand{\Zp}{\mathbb{Z}_p}
\newcommand{\R}{\mathbb{R}}
\newcommand{\C}{\mathbb{C}}
\newcommand{\ii}{\mathrm{i}}
\newcommand{\diag}{\mathrm{diag}}
\newcommand{\E}{\mathbb{E}}

\title{Diagnostic Metrics for the Grokking Phase Transition\\[4pt]
\large Mathematical details, expected signatures, and counterfactual arguments\\
for a Transformer trained on Fibonacci modulo $p$}
\author{Pietro Dall'Olio --- Phase Transitions in Neural Networks project}
\date{May 2026}

\begin{document}
\maketitle
\tableofcontents
\newpage

% ============================================================
\section{Overview of the task and the grokking phenomenon}
% ============================================================

We train a decoder-only Transformer to predict the next token of a Fibonacci
sequence modulo a prime $p$:
\begin{equation}
    x_{t+1} = x_t + x_{t-1} \pmod{p}.
\end{equation}
The vocabulary is $\Zp = \{0,1,\dots,p-1\}$.
Following Nanda et al.\ (2023) we use $p = 113$ and train on a random
30\% of all $p^2 = 12{,}769$ possible consecutive pairs $(a, b)$, evaluating
on the held-out 70\%.

\paragraph{The grokking signature.}
The model typically undergoes two distinct phases:
\begin{enumerate}[label=(\roman*)]
    \item \textbf{Memorisation phase} --- training accuracy quickly reaches 100\%
    while validation accuracy stays near chance $1/p \approx 0.9\%$.
    The model has fit the training pairs without learning the underlying rule.
    \item \textbf{Grokking} --- after many more gradient steps, validation accuracy
    suddenly jumps to $\sim 100\%$.
    The model has discovered the modular-addition structure.
\end{enumerate}
The five families of diagnostic plots described below are designed to detect and
characterise this transition from the inside.

% ============================================================
\section{Loss and accuracy curves}
\label{sec:loss}
% ============================================================

\subsection{Definitions}

We evaluate cross-entropy loss and per-token accuracy, skipping position 0
(where the Fibonacci rule is underdetermined from a single token):
\begin{align}
    \mathcal{L}_{\text{train/val}}
    &= -\frac{1}{|\mathcal{D}|}\sum_{(x,y)\in\mathcal{D}}
       \sum_{t=1}^{T-1} \log \text{softmax}(\ell_\theta(x_{\leq t}))_{y_t},\\
    \text{Acc}_{\text{train/val}}
    &= \frac{1}{|\mathcal{D}|(T-1)}
       \sum_{(x,y)\in\mathcal{D}} \sum_{t=1}^{T-1}
       \mathbf{1}\!\left[\arg\max_c \ell_\theta(c\mid x_{\leq t}) = y_t\right].
\end{align}
Both metrics are plotted on a \emph{log scale} for the loss axis to make the
slow late-stage drop in validation loss visible.

\subsection{Expected behaviour at grokking}

\begin{itemize}
    \item Training loss drops rapidly to near zero; training accuracy hits 100\%.
    \item Validation loss stays near $\log p \approx 4.7$ nats for a long time
    (random-chance prediction).
    \item Then, at the grokking epoch $t_g$, validation loss drops sharply and
    validation accuracy rises to $\sim 100\%$ in a window of a few hundred epochs.
\end{itemize}

\subsection{Counterfactual: pure memorisation, no grokking}

If the model never generalises (e.g.\ weight decay is too small, training is
stopped early, or model capacity is insufficient), the curves have a
characteristic shape:
\begin{itemize}
    \item Train loss converges to $\approx 0$; train accuracy $\approx 100\%$.
    \item Val loss \emph{stays} near $\log p$ indefinitely.
    \item There is no gap closure: the train-val accuracy gap never narrows.
\end{itemize}
Increasing weight decay $\lambda$ is the standard way to force the model to
generalise: it penalises large norms, preventing the model from storing all
training pairs directly in the weights.

% ============================================================
\section{Diagonal spectral mass $M_{\diag}$}
\label{sec:mdiag}
% ============================================================

\subsection{Why the pair lattice $\Zp^2$?}

The Fibonacci rule depends only on the two most recent tokens, so the
sufficient state is the pair $(a,b)\in\Zp^2$.
After grokking the model must have learned a function on this
$p\times p$ torus.
We expose it by building, for every $(a,b)$, a valid probe sequence ending
in $(a,b)$ and collecting the model's final-position logits:
\begin{equation}
    L_\theta[a,b,c] = \ell_\theta(c \mid \text{sequence ending in }(a,b)).
\end{equation}

\subsection{Centered logits and the structure factor}

Softmax is gauge-invariant: adding the same constant to all class logits at
$(a,b)$ does not change predictions.
We therefore center:
\begin{equation}
    \tilde{L}[a,b,c]
    = L[a,b,c] - \frac{1}{p}\sum_{c'} L[a,b,c'].
\end{equation}
The class-summed 2-D power spectrum is
\begin{equation}
    S(k,\ell) = \sum_{c=0}^{p-1}
    \left|\widehat{\tilde{L}}(k,\ell;c)\right|^2,
    \qquad (k,\ell)\in\Zp^2,
\end{equation}
where $\widehat{\cdot}$ denotes the 2-D discrete Fourier transform over
$(a,b)$.

\subsection{Why diagonal modes are special}

The exact target rule is $c = a+b \pmod{p}$.
Its Fourier decomposition is
\begin{equation}
    \mathbf{1}\{c=a+b\} = \frac{1}{p}\sum_{m=0}^{p-1}
    e^{2\pi\ii m(a+b-c)/p}
    = \frac{1}{p}\sum_{m=0}^{p-1}
    e^{-2\pi\ii mc/p}\,
    e^{2\pi\ii ma/p}\,
    e^{2\pi\ii mb/p}.
\end{equation}
The dependence on $(a,b)$ lies entirely on modes $(k,\ell)=(m,m)$ --- the
\emph{diagonal} of the 2-D frequency space.
This is the algebraic reason that a grokked model should concentrate spectral
power on the diagonal.

\subsection{Definition}

\begin{equation}
    M_{\diag}(\theta)
    = \frac{\displaystyle\sum_{m=1}^{p-1} S(m,m)}
           {\displaystyle\sum_{(k,\ell)\neq(0,0)} S(k,\ell)}.
    \label{eq:mdiag}
\end{equation}
$M_{\diag}\in[0,1]$.
The DC mode $(0,0)$ is excluded from both numerator and denominator because it
measures a uniform offset, not spatial structure.

\subsection{Expected behaviour at grokking}

\begin{itemize}
    \item \textbf{Before grokking}: $M_{\diag}\approx 0$.
    The spectrum is spread over all $(k,\ell)$ with no preference for the
    diagonal; the model's output varies irregularly over the pair lattice.
    \item \textbf{At the grokking transition}: $M_{\diag}$ rises sharply,
    tracking the jump in validation accuracy.
    \item \textbf{After grokking}: $M_{\diag}\to 1$.
    Nearly all off-DC spectral power sits on diagonal modes; the model has
    internalised the $a+b \pmod{p}$ structure.
\end{itemize}

\subsection{Counterfactual: memorisation without generalisation}

When the model memorises the training set without learning the rule,
the logit table $L[a,b,c]$ is effectively correct on training pairs and
arbitrary on validation pairs.
The 2-D FFT of a patchy function (correct on a random 30\% of entries,
random on the rest) does \emph{not} concentrate on the diagonal: the
``missing'' 70\% of the pair lattice contributes broad-spectrum noise that
washes out the diagonal structure.
Therefore $M_{\diag}$ remains near a small baseline value ($\sim 0.3 \times
(\text{train fraction})^2$ at best, from spurious alignment of the trained
entries).

The key asymmetry: the diagonal structure in Eq.\,\eqref{eq:mdiag} requires
the \emph{entire} pair lattice to be organised coherently according to the
rule.
Memorisation organises only the training subset, so $M_{\diag}$ stays low.

% ------------------------------------------------------------
\subsection{Derivation: spectral power scales as amplitude squared}
\label{sec:amp2}
% ------------------------------------------------------------

Consider a \emph{single-frequency} approximation to the model's logit table,
parameterised by frequency $f^*$ and (logit) amplitude $A$:
\begin{equation}
    L(a,b,c) \approx A\cos\!\left(\frac{2\pi f^*(a+b-c)}{p}\right).
    \label{eq:single-freq-logit}
\end{equation}
Since the mean over $c$ vanishes (for $p$ prime and $f^*\neq 0$), the
centered logit $\tilde{L}=L$.  Expanding the cosine,
\begin{equation}
    L(a,b,c)
    = \frac{A}{2}\,e^{-2\pi\ii f^* c/p}\,
         e^{2\pi\ii f^* a/p}\,e^{2\pi\ii f^* b/p}
    + \mathrm{c.c.}
\end{equation}
Apply the 2-D DFT over $(a,b)$ with the convention
$\widehat{L}(k,\ell;c) = p^{-1}\sum_{a,b}L(a,b,c)\,
e^{-2\pi\ii(ka+\ell b)/p}$.
Using the discrete orthogonality relation
$p^{-1}\sum_{a}e^{2\pi\ii(f-k)a/p} = \delta_{k,f\bmod p}$, we obtain
\begin{equation}
    \widehat{L}(k,\ell;c)
    = \frac{A}{2}\,e^{-2\pi\ii f^* c/p}\,\delta_{k,f^*}\,\delta_{\ell,f^*}
    + \frac{A}{2}\,e^{+2\pi\ii f^* c/p}\,\delta_{k,p-f^*}\,\delta_{\ell,p-f^*}.
\end{equation}
The structure factor at the diagonal mode $(f^*,f^*)$ is then
\begin{align}
    S(f^*,f^*)
    &= \sum_{c=0}^{p-1}\bigl|\widehat{L}(f^*,f^*;c)\bigr|^2
     = \sum_{c=0}^{p-1}\left(\frac{A}{2}\right)^{\!2}
     = \frac{p\,A^2}{4}.
    \label{eq:Sff}
\end{align}
\textbf{Result:}
\begin{equation}
    \boxed{S(f^*,f^*) = \frac{p}{4}\,A_{f^*}^2.}
    \label{eq:Sscales}
\end{equation}
For a model using $K$ frequencies $\{f^*_i\}$, the cross-terms vanish by
orthogonality of distinct exponentials, so
$S(f^*_i,f^*_i) = (p/4)\,A_{f^*_i}^2$ independently.
The \emph{ratio} of diagonal powers equals the ratio of squared logit
amplitudes:
\begin{equation}
    \frac{S(f^*_1,f^*_1)}{S(f^*_2,f^*_2)} = \frac{A_{f^*_1}^2}{A_{f^*_2}^2}.
\end{equation}

\paragraph{Connection to embedding amplitudes.}
In the idealised grokked circuit (Nanda et al.\ \S4), the attention--MLP block
bilinearly combines embedding vectors and the unembedding reads off the result,
so the logit amplitude scales as $A_f \propto \bigl(A^{(W_E)}_f\bigr)^2$ and
$S(f,f)\propto \bigl(A^{(W_E)}_f\bigr)^4$.
In practice the quartic scaling is only approximate, but the embedding spectrum
$A^{(W_E)}_f$ already gives a reliable ranking of which frequencies are active.

% ------------------------------------------------------------
\subsection{Why the weakest frequency is barely visible in $S(k,k)$}
\label{sec:weak-freq}
% ------------------------------------------------------------

After grokking our model uses five frequencies with embedding amplitudes
approximately
$A_3\approx 2.35,\;A_7\approx 4.2,\;A_{14}\approx 4.0,\;
 A_{18}\approx 4.1,\;A_{56}\approx 4.0$
(indices for $p=113$; exact values are seed-dependent).
By Eq.\,\eqref{eq:Sscales}, the diagonal power ratio of the weakest
to the strongest frequency is
\begin{equation}
    \frac{S(3,3)}{S(7,7)}
    \approx \left(\frac{A_3}{A_7}\right)^{\!2}
    \approx \left(\frac{2.35}{4.2}\right)^{\!2}
    \approx 0.31.
\end{equation}
The $k=3$ peak carries only $\sim 31\%$ of the power of the dominant peaks.

\paragraph{Why this is less severe in the embedding spectrum.}
The embedding spectrum $A(f)$ plots amplitudes linearly; the same ratio is
$2.35/4.2 \approx 0.56$ --- a $56\%$ peak that is clearly visible.
In $S(k,k)$ the ratio is \emph{squared} to $\approx 31\%$, pushing the
$k=3$ bar much closer to the noise floor.  Plotting $S(k,k)$ on a
\emph{log scale} is the standard remedy.

\paragraph{Near-Nyquist merging.}
At $p=113$ the mirror-image modes at $k=56$ and $k=57$ (both near the
Nyquist point $p/2\approx 56.5$) are adjacent on the frequency axis.
If their amplitudes are comparable they visually merge into one broad
feature, further crowding the spectrum and obscuring the individual
contributions.

% ============================================================
\section{Correlation length $\xi$ (Ornstein-Zernike estimator)}
\label{sec:xi}
% ============================================================

\subsection{Intuition}

$M_{\diag}$ tells us \emph{what fraction} of power is on the diagonal.
It does not say \emph{how sharply} that power is concentrated around any
particular mode.
The correlation length $\xi$ captures the sharpness of the dominant spectral
peak, and via the Ornstein-Zernike (OZ) analogy it tells us how spatially
coherent the model's representation of the pair lattice is.

\subsection{The Ornstein-Zernike form}

Near a continuous phase transition the structure factor takes the form
\begin{equation}
    S(k) \approx \frac{A}{\xi^{-2} + |k - k^*|^2},
\end{equation}
where $k^*$ is the dominant ordering wavevector.
In our setting $k^*$ lives on the diagonal.
The peak height scales as $A\xi^2$ and the peak width (half-width at
half-maximum) scales as $\xi^{-1}$.
Therefore:
\begin{itemize}
    \item large $\xi$ $\Rightarrow$ tall, narrow peak $\Rightarrow$
    long-range coherent order on the pair lattice,
    \item small $\xi$ $\Rightarrow$ broad, low peak $\Rightarrow$
    short-range or no order.
\end{itemize}

\subsection{Second-moment estimator on the torus}

On a finite periodic lattice the lattice momenta are
$k_\mu = 2\pi m_\mu / p$ with $m_\mu\in\Zp$.
Using the OZ form and comparing the peak to its four nearest-momentum
neighbours,
\begin{equation}
    \xi^2 \approx
    \frac{1}{4\sin^2(\pi/p)}
    \left(\frac{S(k^*)}{\overline{S}_{\mathrm{nn}}(k^*)} - 1\right),
    \label{eq:xi}
\end{equation}
where $\overline{S}_{\mathrm{nn}}(k^*)$ is the mean of $S$ over the four
nearest neighbours of $k^*$ on the torus.
The factor $4\sin^2(\pi/p)$ replaces the continuum $|\delta k|^2 = (2\pi/p)^2$
with the exact lattice dispersion.

\subsection{Expected behaviour at grokking}

\begin{itemize}
    \item \textbf{Memorisation phase}: $\xi$ is small (order 1 in units of
    the lattice spacing $1/p$) or even undefined if the dominant peak has no
    clear nearest-neighbour structure.
    \item \textbf{Grokking transition}: $\xi$ rises sharply.
    \item \textbf{Post-grokking}: $\xi$ becomes comparable to the system size
    $p$ (the pair lattice is fully coherent); the peak is as sharp as the
    finite lattice allows.
\end{itemize}
Plotting $\xi(t)/p$ for several values of $p$ and looking for a common crossing
is a standard finite-size-scaling test for a continuous transition.

\subsection{Counterfactual: memorisation}

In the memorisation phase the dominant non-DC peak need not even lie on the
diagonal.
Even if it does (because training pairs are a biased sample), the peak is
broad: the off-diagonal noise from untrained pairs competes with it.
Equation \eqref{eq:xi} then gives $S(k^*)/S_{\mathrm{nn}}\approx 1$ and
$\xi\approx 0$.

\paragraph{Joint signature.}
A rising $M_{\diag}$ alone is weak evidence (it could come from a broad
diagonal band).
A simultaneously rising $\xi$ is much stronger: it means the spectral mass
is condensing into a \emph{single narrow} diagonal peak, which is the
hallmark of the model having learned one clean periodic structure.

% ------------------------------------------------------------
\subsection{First-principles derivation of the OZ structure factor}
\label{sec:oz-deriv}
% ------------------------------------------------------------

\paragraph{Ginzburg-Landau setup.}
Define the complex order parameter field
$\phi: \Zp^2 \to \C$,
\begin{equation}
    \phi(a,b) = \sum_{c=0}^{p-1} \tilde{L}(a,b,c)\,e^{-2\pi\ii f^* c/p},
\end{equation}
i.e.\ the Fourier coefficient of the logit table at the target class
frequency $f^*$.
In the grokked phase $\phi(a,b) = A e^{2\pi\ii f^*(a+b)/p}$; in the
disordered phase $\phi$ fluctuates without long-range coherence.

Near the grokking transition (treating it as a continuous phase transition),
expand the free-energy functional in $\phi$ to quadratic order --- the
Ginzburg-Landau functional:
\begin{equation}
    \mathcal{F}[\phi]
    = \sum_{(a,b)\in\Zp^2}
    \Bigl[r\,|\phi(a,b)|^2
         + \kappa\,|\nabla_{\mathrm{lat}}\phi(a,b)|^2
         + \cdots\Bigr],
    \label{eq:GL}
\end{equation}
where $r$ is the (reduced) distance to the transition ($r\propto
\lambda - \lambda_c$ with $\lambda$ the weight decay), and
$\nabla_{\mathrm{lat}}$ denotes the discrete lattice gradient.
The Fourier-space representation of $\mathcal{F}$ is diagonal:
\begin{equation}
    \mathcal{F}[\phi]
    = \sum_{(k,\ell)\in\Zp^2}
    \bigl(r + \kappa\,\mu_{k,\ell}\bigr)\,|\tilde\phi(k,\ell)|^2,
\end{equation}
where $\mu_{k,\ell}$ is the eigenvalue of the lattice Laplacian
$(-\Delta_{\mathrm{lat}})$.

\paragraph{Lattice Laplacian eigenvalue.}
On $\Zp^2$ with periodic boundary conditions the Laplacian acts as
$(-\Delta_{\mathrm{lat}})f = 4f - (\text{four neighbours of }f)$.
Its Fourier eigenvalue at wavevector $(k,\ell)$ is
\begin{equation}
    \mu_{k,\ell}
    = 4 - 2\cos\!\tfrac{2\pi k}{p} - 2\cos\!\tfrac{2\pi\ell}{p}
    = 4\sin^2\!\tfrac{\pi k}{p} + 4\sin^2\!\tfrac{\pi\ell}{p},
    \label{eq:dispersion}
\end{equation}
using $1-\cos\theta = 2\sin^2(\theta/2)$.

\paragraph{Equipartition and the OZ propagator.}
Treating $\mathcal{F}$ as defining a Boltzmann weight
$e^{-\mathcal{F}}$, the Gaussian fluctuation theorem (equipartition) gives
\begin{equation}
    \bigl\langle|\tilde\phi(k,\ell)|^2\bigr\rangle
    = \frac{1}{r + \kappa\,\mu_{k,\ell}}.
\end{equation}
The structure factor $S(k,\ell)\propto
\langle|\tilde\phi(k,\ell)|^2\rangle$ takes the form
\begin{equation}
    S(k,\ell)
    = \frac{A}{r/\kappa + \mu_{k,\ell}}
    = \frac{A}{\xi^{-2}
        + 4\sin^2\!\tfrac{\pi(k-f^*)}{p}
        + 4\sin^2\!\tfrac{\pi(\ell-f^*)}{p}},
    \label{eq:oz-lattice}
\end{equation}
where we shifted to modes relative to the dominant ordering wavevector
$k^*=(\ell^*=f^*)$ and identified the (lattice) \emph{correlation length}
\begin{equation}
    \xi^2 = \kappa/r.
\end{equation}
Equation~\eqref{eq:oz-lattice} is the \textbf{lattice Ornstein-Zernike form}.
In the continuum limit $\sin(\pi\delta k/p)\approx\pi\delta k/p$, so
$\mu\approx|\delta\mathbf{k}|^2$ and we recover the familiar
Lorentzian propagator
\begin{equation}
    S(k) \approx \frac{A}{\xi^{-2} + |\mathbf{k}-\mathbf{k}^*|^2}.
\end{equation}

\paragraph{Non-monotonic $\xi$ profile.}
At the transition $r\to 0^+$, so $\xi\to\infty$ and the peak diverges
(in the infinite-volume, mean-field limit).
On a finite lattice of size $p$, $\xi$ saturates at $O(p)$ once the peak
width reaches the minimum allowed by the Nyquist spacing $2\pi/p$.
After the transition the model settles into the ordered phase with a fixed,
large $\xi$.
The observed non-monotonic profile (rising to a maximum near the grokking
epoch, then gently decreasing) reflects: (i) the finite-size saturation
during the rise; (ii) a slight broadening of the peak once all $K$
frequencies reach comparable amplitudes and
the single-peak approximation becomes less accurate.

% ------------------------------------------------------------
\subsection{Derivation of the second-moment correlation length estimator}
\label{sec:xi-deriv}
% ------------------------------------------------------------

Starting from the lattice OZ form~\eqref{eq:oz-lattice}, evaluate $S$ at
the dominant peak and at its nearest neighbours.

\paragraph{At the peak $(k^*, k^*) = (f^*, f^*)$:}
$\sin^2(0)=0$, so
\begin{equation}
    S_{\mathrm{peak}} \equiv S(k^*,\ell^*)
    = \frac{A}{\xi^{-2}}.
    \label{eq:Speak}
\end{equation}

\paragraph{At a nearest-neighbour $(k^*+1,\, k^*)$:}
The two sine-squared terms contribute $4\sin^2(\pi/p)+0$, so
\begin{equation}
    S_{\mathrm{nn}}
    \equiv S(k^*+1,\,k^*)
    = \frac{A}{\xi^{-2} + 4\sin^2(\pi/p)}.
    \label{eq:Snn}
\end{equation}

\paragraph{Solving for $\xi^2$:}
Dividing Eq.\,\eqref{eq:Speak} by Eq.\,\eqref{eq:Snn},
\begin{equation}
    \frac{S_{\mathrm{peak}}}{S_{\mathrm{nn}}}
    = 1 + \xi^2 \cdot 4\sin^2\!\tfrac{\pi}{p}.
\end{equation}
Rearranging:
\begin{equation}
    \boxed{
    \xi^2 = \frac{1}{4\sin^2(\pi/p)}
            \left(\frac{S_{\mathrm{peak}}}{\overline{S}_{\mathrm{nn}}} - 1\right),
    }
    \label{eq:xi-estimator}
\end{equation}
where in practice $\overline{S}_{\mathrm{nn}}$ is the mean over all four
nearest neighbours $\{(k^*\pm 1, k^*),\, (k^*, k^*\pm 1)\}$ to reduce noise.
This is the estimator implemented in \texttt{analysis.py}.

\paragraph{Lattice vs continuum.}
In the continuum limit $\sin(\pi/p)\approx\pi/p$, so the denominator
approaches $4(\pi/p)^2=(2\pi/p)^2$, the squared momentum-lattice spacing.
The lattice sine factor corrects for the non-linear dispersion
relation~\eqref{eq:dispersion} and is important when $\xi$ is not much
larger than a few lattice spacings (i.e.\ when $\xi \ll p/\pi$).

% ============================================================
\section{Token-embedding Fourier spectrum}
\label{sec:emb}
% ============================================================

\subsection{Definition}

Let $W_E\in\R^{p\times d}$ be the token-embedding matrix (row $k$ is the
embedding of token $k$).
Its 1-D discrete Fourier transform along the token axis is
\begin{equation}
    \hat{W}_E(f) = \frac{1}{\sqrt{p}}\sum_{k=0}^{p-1} W_E[k,\cdot]\,
    e^{-2\pi\ii fk/p},
    \qquad f=0,\dots,\lfloor p/2\rfloor.
\end{equation}
The per-frequency amplitude is the $\ell^2$-norm over embedding dimensions:
\begin{equation}
    A(f) = \left\|\hat{W}_E(f)\right\|_2
    = \left(\sum_{j=1}^{d}|\hat{W}_{E,j}(f)|^2\right)^{1/2}.
\end{equation}

\subsection{What it measures}

$A(f)$ quantifies how much the embedding matrix ``uses'' frequency $f$.
If token $k$ is embedded near $(\cos(2\pi f k/p),\,\sin(2\pi f k/p),\dots)$
for some direction in $\R^d$, then $A(f)$ will be large.
A flat spectrum $A(f)\approx\text{const}$ means tokens are embedded without
any periodic structure.

\subsection{Expected behaviour at grokking}

Nanda et al.\ show that after grokking the embedding concentrates at a small
number of \emph{key frequencies} $f^*$.
Concretely, a grokked model embeds token $k$ as
\begin{equation}
    W_E[k,:] \approx \sum_{f\in\mathcal{F}} \alpha_f \cos(2\pi f k/p)\,u_f
    + \beta_f \sin(2\pi f k/p)\,v_f,
\end{equation}
where $\mathcal{F}$ is a small set (typically 2--5 frequencies),
$u_f,v_f\in\R^d$ are ``directions'' learned by the model, and
$\alpha_f,\beta_f$ are scalar amplitudes.
The Fourier spectrum $A(f)$ therefore shows a few sharp peaks above a low
noise floor.

\begin{itemize}
    \item \textbf{Before grokking}: $A(f)$ is nearly flat (no preferred
    frequencies; embeddings are arbitrary high-dimensional vectors adjusted
    to memorise specific pairs).
    \item \textbf{After grokking}: $A(f)$ has a small number of dominant
    spikes.
    The positions $f^*$ of these spikes are the same frequencies that
    dominate the diagonal of $S(k,\ell)$.
\end{itemize}

\subsection{Counterfactual: memorisation}

A model that has only memorised the training set does not need a periodic
embedding: it just needs token $k$ to have a distinct representation.
The embedding will therefore be spread across many Fourier modes (flat
spectrum), reflecting no shared structure between tokens beyond their
distinctness.

\subsection{Relationship to $M_{\diag}$}

The key frequencies $f^*$ appearing in the embedding are precisely the
frequencies $m$ for which $S(m,m)$ is largest.
Both metrics are measuring the same underlying phenomenon from different
angles: $M_{\diag}$ asks ``does the output logit table use diagonal Fourier
modes?'' while $A(f)$ asks ``does the input embedding use the corresponding
Fourier basis?''

% ============================================================
\section{Trigonometric circle representation}
\label{sec:trig}
% ============================================================

\subsection{The prediction mechanism (Nanda et al., §3)}

Once the model has selected key frequencies $\mathcal{F}$, the prediction of
$c = a+b\pmod{p}$ can be implemented almost linearly:
\begin{align}
    \cos(2\pi f c / p)
    &= \cos(2\pi f a/p)\cos(2\pi f b/p) - \sin(2\pi f a/p)\sin(2\pi f b/p),\\
    \sin(2\pi f c / p)
    &= \sin(2\pi f a/p)\cos(2\pi f b/p) + \cos(2\pi f a/p)\sin(2\pi f b/p).
\end{align}
The attention and MLP layers effectively compute these trigonometric products
and then the unembedding reads off the token whose trig components match the
computed result.

\subsection{The circle plot}

For a given frequency $f$ we project $W_E[k,:]$ onto the cosine and sine
basis vectors:
\begin{align}
    \phi^{(c)}_f[k] &= W_E[k,:]\cdot \hat{c}_f, \\
    \phi^{(s)}_f[k] &= W_E[k,:]\cdot \hat{s}_f,
\end{align}
where $\hat{c}_f,\hat{s}_f\in\R^d$ are the normalised cosine and sine
directions at frequency $f$ extracted from $\hat{W}_E(f)$.
The scatter plot of $(\phi^{(c)}_f[k],\,\phi^{(s)}_f[k])$ for all
$k=0,\dots,p-1$, coloured by $k$, should look like an \emph{evenly-spaced
ring} if the model has internalised the trig structure.

\subsection{Expected behaviour at grokking}

\begin{itemize}
    \item \textbf{Before grokking}: the scatter plot is a diffuse cloud with
    no clear circular structure; points are arranged approximately randomly.
    \item \textbf{After grokking}: points lie on a circle and are evenly
    spaced modulo $p$ (i.e.\ token $k$ is at angle $2\pi fk/p$), with
    colours cycling uniformly around the ring.
    The radius of the ring is $|\alpha_f|$ (the amplitude of that frequency
    in the embedding).
\end{itemize}

\subsection{Counterfactual: memorisation}

A memorising model needs to distinguish all $p$ tokens but has no reason
to arrange them periodically.
The circle plot will show a diffuse cloud, or at best a ring whose angular
positions are irregular (tokens are not at $2\pi fk/p$).
Even if there happens to be some residual circular structure (perhaps from
the random initialisation), the coloured sequence of tokens will \emph{not}
cycle uniformly around the ring; adjacent-colour tokens will be scattered.

% ============================================================
\section{Joint interpretation: memorisation vs.\ grokking trajectory}
\label{sec:joint}
% ============================================================

Table~\ref{tab:summary} summarises the expected state of each diagnostic
at three stages of training.

\begin{table}[h]
\centering
\begin{tabular}{lp{3.2cm}p{3.2cm}p{3.8cm}}
\toprule
Metric & Early (random) & Memorisation & Grokking (generalised) \\
\midrule
Train acc & $\approx 1/p$ & $\approx 100\%$ & $\approx 100\%$ \\
Val acc   & $\approx 1/p$ & $\approx 1/p$  & $\approx 100\%$ \\
$M_{\diag}$ & $\approx 0$ & small, noisy & rises sharply $\to 1$ \\
$\xi$       & $\approx 0$ & $\approx 0$--1 & rises to $O(p)$ \\
$A(f)$ spectrum & flat & flat & sharp peaks at $f^*$ \\
Circle plot & diffuse cloud & diffuse / irregular ring & evenly-spaced ring \\
\bottomrule
\end{tabular}
\caption{Summary of diagnostic behaviour at each stage.}
\label{tab:summary}
\end{table}

\subsection{Why memorisation does not produce the spectral signatures}

The core asymmetry is that \emph{the spectral metrics are computed over all
$p^2$ states}, while the model has only seen 30\% of them during training.

\begin{itemize}
    \item The 2-D FFT over the full pair lattice requires all $p^2$ entries
    to be organised coherently for $M_{\diag}$ to be close to 1.
    A model that has fit the training pairs but outputs arbitrary values on
    validation pairs produces a logit table $L[a,b,c]$ that is structured
    on 30\% of the grid and noisy on 70\%.
    The Fourier transform of such a function is dominated by the noise, and
    $M_{\diag}$ remains small.

    \item Similarly, $\xi$ measures the \emph{peak-to-neighbourhood ratio}
    of $S(k,\ell)$.
    Random off-diagonal power elevates the neighbourhood values, keeping the
    ratio near 1 and $\xi$ near 0.

    \item The embedding spectrum $A(f)$ and the circle plot are determined
    by the token embedding $W_E$, which is trained on the full vocabulary
    but must serve as input to the whole model.
    During memorisation, $W_E$ is tuned to produce correct outputs on training
    pairs; it is under no pressure to adopt a globally periodic structure.
\end{itemize}

\subsection{Why rising $M_{\diag}$ precedes or coincides with the val-acc jump}

During training the model implicitly competes between two strategies:
\begin{enumerate}[label=(\alph*)]
    \item Store training pairs as a lookup table (memorisation).
    \item Compress the task into the modular-addition rule using trigonometric
    representations (generalisation).
\end{enumerate}
The spectral metrics detect strategy (b) directly.
Weight decay penalises the large norms required to memorise many individual
pairs, creating an implicit pressure toward the more weight-efficient
(b) strategy.
Once the model commits to (b), the internal representations reorganise
simultaneously, causing $M_{\diag}$, $\xi$, the embedding spectrum, and
val accuracy to jump together within a short training window.

\subsection{What a slow $M_{\diag}$ rise means}

If $M_{\diag}$ rises gradually rather than sharply, the model may be:
\begin{enumerate}
    \item transitioning slowly (crossover rather than sharp transition),
    \item using multiple competing frequencies simultaneously
    (slow ``frequency selection''),
    \item partially generalising via a subset of the diagonal modes.
\end{enumerate}
Checking the structure-factor heatmap and the circle plots at intermediate
checkpoints can distinguish these cases.

% ============================================================
\section{Practical monitoring recipe}
\label{sec:recipe}
% ============================================================

During training, call \texttt{monitor\_spectral\_observables} every
\texttt{monitor\_every} epochs.
Record and plot:
\begin{enumerate}
    \item Training and validation loss / accuracy on log scale.
    \item $M_{\diag}(t)$ and $\xi(t)$ vs.\ epoch.
    \item The structure-factor heatmap $S(k,\ell)$ at selected checkpoints
    (before, during, and after the suspected transition).
    \item The embedding Fourier spectrum $A(f)$ at the same checkpoints.
    \item Circle plots for the top-3 frequencies after training.
\end{enumerate}

\paragraph{Checkpoint selection heuristic.}
The grokking epoch $t_g$ is usually where $dM_{\diag}/dt$ is maximum.
Save full snapshots at $t = 0$, $t_g/2$, $t_g$, and $2t_g$.

\paragraph{Finite-size scaling extension.}
To test whether the transition sharpens with system size, repeat the
experiment for several values of $p$ (e.g.\ $p\in\{17,31,59,97,113\}$)
with the same training fraction and hyperparameters.
Plot $\xi(t)/p$ for all $p$ on the same graph; a common crossing near $t_g$
is evidence of a continuous transition with well-defined critical behaviour.

% ============================================================
\section{Weight decay: role, formulas, and effect of batch size}
\label{sec:weightdecay}
% ============================================================

\subsection{Origin: L2 regularisation}

The motivation for weight decay is to prevent overfitting by penalising large
parameters.  One adds a regularisation term to the loss:
\begin{equation}
    \mathcal{L}_{\mathrm{reg}}(\theta)
    = \mathcal{L}(\theta) + \frac{\lambda}{2}\|\theta\|^2.
\end{equation}
The gradient becomes $\nabla\mathcal{L}_{\mathrm{reg}} = \nabla\mathcal{L} +
\lambda\theta$, and a plain SGD step gives
\begin{equation}
    \theta_{t+1}
    = \theta_t - \eta\,\nabla\mathcal{L}_{\mathrm{reg}}
    = (1-\eta\lambda)\,\theta_t - \eta\,\nabla\mathcal{L}.
    \label{eq:sgd-wd}
\end{equation}
The factor $(1-\eta\lambda)$ is the \emph{per-step weight shrinkage}: every
parameter is multiplied by a number slightly less than 1 before the gradient
correction is added.  In the absence of gradient signal the weights decay
exponentially toward zero --- hence the name \emph{weight decay}.

\subsection{AdamW: decoupled weight decay}

Standard Adam with L2 regularisation applies the penalty \emph{through} the
adaptive gradient scaling, which corrupts the regularisation effect.
AdamW~\citep{loshchilov2019} fixes this by applying weight decay directly to
the parameters, bypassing Adam's $1/\sqrt{\hat{v}}$ rescaling.

Let $g_t = \nabla_\theta\mathcal{L}$ be the loss gradient at step $t$.
Define the first and second gradient moments as
\begin{align}
    m_t &= \beta_1 m_{t-1} + (1-\beta_1)\,g_t, \\
    v_t &= \beta_2 v_{t-1} + (1-\beta_2)\,g_t^2,
\end{align}
with bias-corrected estimates
$\hat{m}_t = m_t/(1-\beta_1^t)$ and $\hat{v}_t = v_t/(1-\beta_2^t)$.
The AdamW parameter update is then
\begin{equation}
    \boxed{\theta_{t+1}
    = (1-\eta\lambda)\,\theta_t
    \;-\; \eta\,\frac{\hat{m}_t}{\sqrt{\hat{v}_t}+\varepsilon}.}
    \label{eq:adamw}
\end{equation}
The weight-decay term $\eta\lambda\theta_t$ is additive and decoupled from
the adaptive part: it shrinks every parameter by the same factor
$(1-\eta\lambda)$ at every step, regardless of gradient magnitude or history.
In PyTorch \texttt{AdamW} the parameter \texttt{weight\_decay} is exactly
$\lambda$.

\subsection{Per-epoch shrinkage and why batch size matters}

Because weight decay is applied at \emph{every gradient step}, the number of
steps per epoch is crucial.  With $N_{\mathrm{train}}$ training samples and
batch size $B$, there are
\begin{equation}
    n_{\mathrm{steps}} = \left\lceil \frac{N_{\mathrm{train}}}{B} \right\rceil
    \quad \text{steps per epoch.}
\end{equation}
After one epoch, the multiplicative shrinkage from weight decay alone ---
ignoring the gradient term --- is
\begin{equation}
    \alpha_{\mathrm{epoch}}
    = (1-\eta\lambda)^{n_{\mathrm{steps}}}
    \approx e^{-\eta\lambda\, n_{\mathrm{steps}}},
    \label{eq:epoch-shrink}
\end{equation}
where the approximation uses $\ln(1-x)\approx -x$ for small $x=\eta\lambda$.

For our defaults ($\eta=10^{-3}$, $\lambda=0.1$, $N_{\mathrm{train}}\approx
3828$, so $\eta\lambda=10^{-4}$):

\begin{center}
\begin{tabular}{lrrr}
\toprule
Batch size $B$ & Steps/epoch & $\alpha_{\mathrm{epoch}}$ & Decay per epoch \\
\midrule
3828 (full batch) & 1   & $0.9999^{1}   \approx 0.9999$ & $0.01\%$ \\
1000              & 4   & $0.9999^{4}   \approx 0.9996$ & $0.04\%$ \\
32                & 120 & $0.9999^{120} \approx 0.988$  & $1.2\%$  \\
\bottomrule
\end{tabular}
\end{center}

With $B=32$ the weights are shrunk 30 times more aggressively per epoch than
with $B=1000$, and 120 times more than with full-batch training.

\subsection{Why this causes oscillations in the loss curve}

The grokked solution lives at a specific point in weight space with specific
trigonometric amplitudes in the embedding and MLP weights
(see Section~\ref{sec:trig}).  Near that point the training loss is
essentially zero, so the gradient is very small.  But weight decay is
indifferent to the loss: it still shrinks every parameter by the same amount.

The competition at each step is between two forces:
\begin{equation}
    \underbrace{\eta\,\frac{\hat{m}_t}{\sqrt{\hat{v}_t}+\varepsilon}}
    _{\text{pulls toward grokked solution}}
    \quad \text{vs.} \quad
    \underbrace{\eta\lambda\,\theta_t}_{\text{pulls toward }0}.
\end{equation}
Near a loss minimum the gradient signal is weak.  With $B=32$ and 120
steps per epoch, weight decay repeatedly wins: the trigonometric amplitudes
erode, the model ``un-grokks'', and the loss spikes.  A large gradient then
drives the model back toward the grokked solution, the loss crashes, and the
cycle repeats.  This is the limit cycle visible in the loss curves when
training with small batch size.

With full-batch training the gradient at every step is computed from all
$N_{\mathrm{train}}$ examples.  Even near the loss minimum it is large and
unbiased enough to counteract weight decay at each step, so the grokked
solution is stable and the curves are smooth.

\subsection{Physical analogy}

Think of weight decay as a restoring force (a spring) pulling $\theta$
toward $0$, and the gradient descent step as a driving force toward the loss
minimum.  The grokked fixed point is stable only when the driving force is
stronger than the spring.  With mini-batches the driving force is noisy and
weak near the minimum, so the spring intermittently wins and the system
oscillates.

This analogy also illuminates why weight decay is \emph{necessary} for
grokking, not just harmful.  Without weight decay the memorisation basin
(high-norm weights that encode each training pair individually) has no penalty
and the model stays there indefinitely.  Weight decay provides the pressure
that eventually forces a transition to the lower-norm grokked solution.  The
challenge is tuning $\lambda$ and $B$ so that the spring is strong enough to
favour the grokked solution over the memorised one, but not so strong that it
destabilises the grokked solution once found.

\subsection{Practical fixes}

\paragraph{Full-batch training.}
Set \texttt{batch\_size} $= N_{\mathrm{train}}$ (approximately 3828 for our
default settings).  This is what Nanda et al.\ use and it completely
eliminates the oscillations.

\paragraph{Rescale weight decay with batch size.}
To keep $\alpha_{\mathrm{epoch}}$ constant across batch sizes, set
\begin{equation}
    \lambda_{\mathrm{eff}} = \lambda \cdot \frac{B}{N_{\mathrm{train}}},
\end{equation}
so that the per-epoch shrinkage is $(1-\eta\lambda_{\mathrm{eff}})^{N_{\mathrm{train}}/B}
= (1-\eta\lambda)^1$, matching the full-batch regime.

\paragraph{Apply weight decay once per epoch.}
One can decouple weight decay from the gradient step entirely by applying
the shrinkage $\theta \leftarrow (1-\eta\lambda)\,\theta$ only once at the
end of each epoch, regardless of how many mini-batches were processed.
This is not standard in PyTorch but is straightforward to implement manually.

\begin{thebibliography}{9}
\bibitem{loshchilov2019}
I.~Loshchilov and F.~Hutter,
\textit{Decoupled Weight Decay Regularization},
ICLR 2019.
\bibitem{nanda2023}
N.~Nanda, L.~Chan, T.~Lieberum, J.~Smith, J.~Steinhardt,
\textit{Progress Measures for Grokking via Mechanistic Interpretability},
ICLR 2023.
\end{thebibliography}

\end{document}
